\section{Discussion}\label{sec:discussion}

The frozen features keep the distinctions that some class boundary demanded in pretraining and lose part of the rest.
The separation $S$ makes this quantitative: label sets that keep apart the classes the \Kfull-class probe confuses keep
the accuracy, whatever their size, and one ordering covers the nested sets, the random groupings and the
flavor-boundary groupings.
A random grouping that merges $X\to b\bar b$ with $X\to c\bar c$ still separates them almost as well as one that keeps
them apart, while the \Kcoarse-class set, which merges bottom with charm in every family, does not.
This suggests that a network learns a flavor distinction wherever any class boundary requires it and then applies it
to decays whose own labels do not.
The residual of section~\ref{sec:which} shows that this is not the whole story: how the classes are grouped matters
beyond the separation the grouping keeps, since the \Kmid-class set falls below groupings with less separation.

The same pattern has been seen in image classification, where coarser labels cost more for frozen features than after
fine-tuning~\cite{hong}, and classes merged along a semantic hierarchy remain more entangled than classes merged at
random~\cite{yang}.
What simulation adds is control: the labels are exact, the training examples and their order are identical for every
label set, and each merge removes a distinction that can be named, so the loss can be traced to individual pairs of
classes instead of being inferred from an average over a taxonomy.

For an analysis that reads a pretrained tagger out from its frozen features, or sums class scores into a
discriminant, the classes it will need to tell apart should be distinct somewhere in the pretraining
labels~\cite{mahajan}.
Fine-tuning relaxes this requirement without removing it, and less so when labeled jets are few.
Tasks whose inputs lack the flavor information show no dependence on the label set, even where, as for quark--gluon
discrimination, they draw on a distinction that the coarse labels merge.
Random groupings are not a recipe for label design, however: their frozen features transfer to top tagging less well
than those of any nested set and no better to JetClass, which may reflect that their classes mix decays of different
structure.
The mass regression shows that what the labels lack can be supplied by an auxiliary target, at a small cost in
fine-tuned accuracy.

Our study has limits.
It rests on one network, with \VParamsMillion~million parameters, and on one simulated dataset.
Each step of the hierarchy removes several distinctions at once, so the cost of a single distinction is established
only through pairs of classes and the flavor-boundary groupings, and each random grouping was pretrained once.
All fine-tuned comparisons use one recipe, which leaves the \Kfull-class model slightly below its own linear probe on
JetClass-II; we have not measured whether a recipe that keeps more of the pretrained output layer changes the gaps
between label sets.
On top tagging the same recipe stays below the best published rejections (appendix~\ref{app:published}), so the
absolute performance on the benchmarks should not be read as the reach of the method.

\section{Conclusion}\label{sec:conclusion}

Varying only the labels of a supervised pretraining shows that the representation is shaped by which distinctions the
labels require, and much less by how many classes there are.
A label set that separates the classes later analyses need, anywhere in its structure, keeps them separable in the
frozen features, and one that merges them loses part of that separation, which fine-tuning on limited labeled data
recovers only in part.
For jet tagging, where simulation makes fine labels cheap, pretraining labels should keep apart the flavors and decay
products that analyses will select on, above all when the model is to be used frozen, and quantities that no label
carries, such as the jet mass, can be added as auxiliary targets.
